% GENERATED FILE. Edit canonical lessons, then run npm run generate:books.
% canonical-source-hash: fnv1a64:47d4cdd166dfca94
% canonical-lessons: LA-C347-monstrum, LA-C347-gigas, LA-C347-heros, LA-C347-nympha, LA-C347-vellus

\chapter{Mōnstra et hērōēs}
\label{ch:la-monstra-et-heroes}
\hlchaptermodality{347}

\begin{chapteropening}
\textbf{By the end of this chapter:} \emph{I can talk about the creatures of myth: a monster, a giant, a hero, a nymph and the fleece.}

Complete the last lesson of chapter 347: I can talk about the creatures of myth: a monster, a giant, a hero, a nymph and the fleece.
\end{chapteropening}

\section[mōnstrum, mōnstrī]{mōnstrum, mōnstrī — a monster; a portent}

\label{lesson:LA-C347-monstrum}



\begin{quote}
Before the new one: say the Latin for mortal, then the Latin for immortal.
\end{quote}

\subsection*{You'll want to know: mōnstrum, mōnstrī}
\textbf{mōnstrum, mōnstrī} — \textquotedblleft{}a monster; a portent\textquotedblright{}.

\begin{quote}
Mōnstrum ingēns in spēluncā habitābat.
\end{quote}

\begin{quote}
\emph{A huge monster lived in the cave.}
\end{quote}

Neuter, from \emph{moneō}, 'to warn': first a sign sent by the gods, then a monster. English \emph{monster}.

\begin{cousinweb}
\textbf{mōnstrum} is built on \emph{moneō}, \textquotedblleft{}to warn\textquotedblright{} (EtymDB, Lewis \& Short): a warning sign, a portent, and then a monster. English \emph{monster} comes from it through French.
\end{cousinweb}

\begin{grammarlens}[title={What you've built}]
One more word for telling, explaining and describing at B1.
\end{grammarlens}

\subsection*{Your turn}
\begin{itemize}
\raggedright
  \item \emph{Say it:} \emph{mōnstrum, mōnstrī}
  \item \emph{Say it:} \emph{mōnstrum, mōnstrī}, once more
  \item \emph{Say it:} a sentence of your own with \emph{mōnstrum, mōnstrī}
  \item \emph{Recall:} say the Latin for mortal, then the Latin for immortal, then say \emph{mōnstrum, mōnstrī} again
\end{itemize}

\begin{tcolorbox}[breakable,colback=teal!4,colframe=teal!35!black,title={Before you move on}]
What does mōnstrum, mōnstrī mean? (\textquotedblleft{}A monster; a portent\textquotedblright{}.) Say it once more.
\end{tcolorbox}
\section[gigās, gigantis]{gigās, gigantis — a giant}

\label{lesson:LA-C347-gigas}



\begin{quote}
Before the new one: say the Latin for immortal, then the Latin for a monster.
\end{quote}

\subsection*{You'll want to know: gigās, gigantis}
\textbf{gigās, gigantis} — \textquotedblleft{}a giant\textquotedblright{}.

\begin{quote}
Gigantēs saxa ad caelum iaciēbant.
\end{quote}

\begin{quote}
\emph{The giants were hurling rocks at the sky.}
\end{quote}

Masculine, a Greek loan; usually plural. English \emph{giant}.

\begin{cousinweb}
\textbf{gigās} was borrowed from Greek \emph{gigas}, \textquotedblleft{}giant\textquotedblright{} (EtymDB). English \emph{giant} comes from it through French.
\end{cousinweb}

\begin{grammarlens}[title={What you've built}]
One more word for telling, explaining and describing at B1.
\end{grammarlens}

\subsection*{Your turn}
\begin{itemize}
\raggedright
  \item \emph{Say it:} \emph{gigās, gigantis}
  \item \emph{Say it:} \emph{gigās, gigantis}, once more
  \item \emph{Say it:} a sentence of your own with \emph{gigās, gigantis}
  \item \emph{Recall:} say the Latin for immortal, then the Latin for a monster, then say \emph{gigās, gigantis} again
\end{itemize}

\begin{tcolorbox}[breakable,colback=teal!4,colframe=teal!35!black,title={Before you move on}]
What does gigās, gigantis mean? (\textquotedblleft{}A giant\textquotedblright{}.) Say it once more.
\end{tcolorbox}
\section[hērōs, hērōis]{hērōs, hērōis — a hero, a demigod}

\label{lesson:LA-C347-heros}



\begin{quote}
Before the new one: say the Latin for a monster, then the Latin for a giant.
\end{quote}

\subsection*{You'll want to know: hērōs, hērōis}
\textbf{hērōs, hērōis} — \textquotedblleft{}a hero, a demigod\textquotedblright{}.

\begin{quote}
Hērōs mōnstrum gladiō interfēcit.
\end{quote}

\begin{quote}
\emph{The hero killed the monster with a sword.}
\end{quote}

Masculine, a Greek loan; accusative \emph{hērōa}. A man of the old stories, often the son of a god.

\begin{cousinweb}
\textbf{hērōs} was borrowed from Greek \emph{hērōs}, \textquotedblleft{}demigod, hero\textquotedblright{} (EtymDB). English \emph{hero} comes from it.
\end{cousinweb}

\begin{grammarlens}[title={What you've built}]
One more word for telling, explaining and describing at B1.
\end{grammarlens}

\subsection*{Your turn}
\begin{itemize}
\raggedright
  \item \emph{Say it:} \emph{hērōs, hērōis}
  \item \emph{Say it:} \emph{hērōs, hērōis}, once more
  \item \emph{Say it:} a sentence of your own with \emph{hērōs, hērōis}
  \item \emph{Recall:} say the Latin for a monster, then the Latin for a giant, then say \emph{hērōs, hērōis} again
\end{itemize}

\begin{tcolorbox}[breakable,colback=teal!4,colframe=teal!35!black,title={Before you move on}]
What does hērōs, hērōis mean? (\textquotedblleft{}A hero, a demigod\textquotedblright{}.) Say it once more.
\end{tcolorbox}
\section[nympha, nymphae]{nympha, nymphae — a nymph}

\label{lesson:LA-C347-nympha}



\begin{quote}
Before the new one: say the Latin for a giant, then the Latin for a hero.
\end{quote}

\subsection*{You'll want to know: nympha, nymphae}
\textbf{nympha, nymphae} — \textquotedblleft{}a nymph\textquotedblright{}.

\begin{quote}
Nymphae in silvīs et fontibus habitābant.
\end{quote}

\begin{quote}
\emph{Nymphs lived in the woods and springs.}
\end{quote}

Feminine, a Greek loan; also 'a bride, a young woman'. Spirits of springs, trees and mountains.

\begin{cousinweb}
\textbf{nympha} was borrowed from Greek \emph{nymphē}, \textquotedblleft{}bride, nymph\textquotedblright{} (Lewis \& Short). English \emph{nymph} comes from it.
\end{cousinweb}

\begin{grammarlens}[title={What you've built}]
One more word for telling, explaining and describing at B1.
\end{grammarlens}

\subsection*{Your turn}
\begin{itemize}
\raggedright
  \item \emph{Say it:} \emph{nympha, nymphae}
  \item \emph{Say it:} \emph{nympha, nymphae}, once more
  \item \emph{Say it:} a sentence of your own with \emph{nympha, nymphae}
  \item \emph{Recall:} say the Latin for a giant, then the Latin for a hero, then say \emph{nympha, nymphae} again
\end{itemize}

\begin{tcolorbox}[breakable,colback=teal!4,colframe=teal!35!black,title={Before you move on}]
What does nympha, nymphae mean? (\textquotedblleft{}A nymph\textquotedblright{}.) Say it once more.
\end{tcolorbox}
\section[vellus, velleris]{vellus, velleris — a fleece}

\label{lesson:LA-C347-vellus}



\begin{quote}
Before the new one: say the Latin for a hero, then the Latin for a nymph.
\end{quote}

\subsection*{You'll want to know: vellus, velleris}
\textbf{vellus, velleris} — \textquotedblleft{}a fleece\textquotedblright{}.

\begin{quote}
Hērōēs vellus aureum per tōtum mare quaesīvērunt.
\end{quote}

\begin{quote}
\emph{The heroes searched the whole sea for the golden fleece.}
\end{quote}

Neuter, 3rd declension. A sheep's shorn wool; the golden one was guarded by a serpent.

\begin{cousinweb}
\textbf{vellus} is inherited from a Proto-Indo-European word for \textquotedblleft{}wool, hair\textquotedblright{} (EtymDB); English \emph{wool} is a cousin.
\end{cousinweb}

\begin{grammarlens}[title={What you've built}]
One more word for telling, explaining and describing at B1.
\end{grammarlens}

\subsection*{Your turn}
\begin{itemize}
\raggedright
  \item \emph{Say it:} \emph{vellus, velleris}
  \item \emph{Say it:} \emph{vellus, velleris}, once more
  \item \emph{Say it:} a sentence of your own with \emph{vellus, velleris}
  \item \emph{Recall:} say the Latin for a hero, then the Latin for a nymph, then say \emph{vellus, velleris} again
\end{itemize}

\begin{tcolorbox}[breakable,colback=teal!4,colframe=teal!35!black,title={Before you move on}]
What does vellus, velleris mean? (\textquotedblleft{}A fleece\textquotedblright{}.) Say it once more.
\end{tcolorbox}
